\section{Conceptual Foundations: Information Disorder and Synthetic Personas}
\label{sec:conceptual_foundations}

To discuss the impact of LLMs on media ecosystems, this paper uses the framework of \emph{information disorder} by Wardle and Derakhshan~\cite{wardle2017information}, written for the Council of Europe. It distinguishes three types of problematic content along two questions, whether the content is false and whether it is shared with the intent to harm:
\begin{itemize}
    \item \textbf{Misinformation:} false content shared without the intent to cause harm. For LLMs, unintentional factual errors (``hallucinations'') fall into this category.
    \item \textbf{Disinformation:} false content that is created or shared deliberately to deceive or cause harm, for example fabricated news articles.
    \item \textbf{Malinformation:} genuine content that is used to cause harm, for example by moving it out of its original context.
\end{itemize}
Organised, often automated, efforts to manipulate public opinion through social media are studied under the label \emph{computational propaganda}~\cite{woolley2018computational}. Earlier work already argued that language models could make the writing of such content cheaper and easier to scale~\cite{buchanan2021truth}.

In parallel, LLMs are studied as \emph{synthetic personas}: models are prompted to answer as a person or as a member of a group, and their answers are compared with human data~\cite{aher2023using, santurkar2023whose}. Whether such answers reflect the diversity of real human populations, or instead the preferences built into the model during training, is the common thread of Sections~\ref{sec:persona_simulation} and~\ref{sec:demographic_alignment}. Figure~\ref{fig:framework} summarises how the four reviewed studies relate to these concepts.

% Figure 1: conceptual framework of the review. All statements in the study boxes are taken from the
% reviewed papers (page references in the section texts); the two yellow boxes are this paper's own synthesis.
\begin{figure}[t]
\centering
\begin{tikzpicture}[
  font=\footnotesize,
  root/.style={draw=black!70, fill=black!8, rounded corners=2pt, align=center, inner sep=4pt, text width=6.0cm},
  role/.style={draw=#1!70!black, fill=#1!12, rounded corners=2pt, align=center, inner sep=3pt, text width=7.0cm, font=\footnotesize\bfseries},
  study/.style={draw=#1!60!black, fill=white, rounded corners=2pt, align=left, inner sep=3pt, text width=3.2cm, anchor=north},
  synth/.style={draw=black!70, fill=yellow!15, rounded corners=2pt, align=center, inner sep=3pt, text width=7.0cm, anchor=north},
  arr/.style={-{Latex[length=2mm]}, black!60, thick}
]
\node[root] (llm) at (0,0) {\textbf{Instruction-tuned / RLHF LLMs}\\ fluent, steerable text generation~\cite{ouyang2022training}};
\node[role=blue] (r1) at (-3.95,-1.35) {Role 1: LLMs as synthetic personas};
\node[role=red]  (r2) at ( 3.95,-1.35) {Role 2: LLMs and information disorder};

\node[study=blue] (aher) at (-5.9,-1.85) {\textbf{Aher et al.}~\cite{aher2023using}\\ Turing Experiments: 3 of 4 classic findings replicated; \emph{hyper-accuracy distortion} in wisdom-of-crowds task};
\node[study=blue] (sant) at (-2.0,-1.85) {\textbf{Santurkar et al.}~\cite{santurkar2023whose}\\ OpinionQA: LM opinions misaligned with US groups; steering helps only modestly};
\node[study=red]  (vyk)  at ( 2.0,-1.85) {\textbf{Vykopal et al.}~\cite{vykopal2024disinformation}\\ Generation: LLMs write convincing disinformation; best detectors $\sim$0.8 macro-F1};
\node[study=red]  (gae)  at ( 5.9,-1.85) {\textbf{Gaeta et al.}~\cite{gaeta2025towards}, SemEval~\cite{dasanmartino2020semeval}\\ Detection: prompted LLMs find propaganda techniques and spans};

\draw[arr] (llm.south) -- ++(0,-0.25) -| (r1.north);
\draw[arr] (llm.south) -- ++(0,-0.25) -| (r2.north);

\node[synth] (ctrl) at (-3.95,-5.0) {\textbf{Controllability gap} (Sec.~\ref{sec:comparative_synthesis})\\ prompting steers analytic tasks more easily than opinion distributions};
\node[synth] (asym) at ( 3.95,-5.0) {\textbf{Generation--detection asymmetry} (Sec.~\ref{sec:comparative_synthesis})\\ cheap generation vs.\ costly, imperfect detection};

\draw[arr, dashed] (aher.south) -- (aher.south |- ctrl.north);
\draw[arr, dashed] (sant.south) -- (sant.south |- ctrl.north);
\draw[arr, dashed] (vyk.south) -- (vyk.south |- asym.north);
\draw[arr, dashed] (gae.south) -- (gae.south |- asym.north);
\end{tikzpicture}
\caption{Conceptual framework of this review. The four studies are grouped by the role of the LLM; the two yellow boxes are the cross-study tensions derived in this paper (own synthesis, not a finding of any single study).}
\label{fig:framework}
\end{figure}

